%% Q1
In \textbf{synchronous} communication all participants are present \emph{at the same time} and interact in real time. Video conferencing (Zoom, Google Meet, Microsoft Teams), live chat, telephone calls and live webinars are synchronous.
\why{E-mails, blogs and wikis are \textbf{asynchronous}: the sender writes now and the receiver reads or edits later, at a time of their own choosing.}
\tip{Synchronous = same time (live). Asynchronous = different time (stored and forwarded).}

%% Q2
\textbf{John McCarthy} (1927--2011) coined the term \emph{Artificial Intelligence} in the 1955 proposal for the \textbf{Dartmouth Summer Research Project (1956)}, the event regarded as the birth of AI as a field. He also created the \textbf{LISP} programming language (1958), long the main language of AI research.
\why{Charles Babbage designed the Analytical Engine and is called the ``father of the computer''. Andrew Ng is a modern machine-learning researcher (co-founder of Coursera and Google Brain). Mark Zuckerberg founded Facebook.}

%% Q3
The 40 SWAYAM Prabha DTH channels are each dedicated to a theme. Channel~01 \emph{Vageesh} (Language \& Literature), Channel~03 \emph{Prabodh} (Social \& Behavioural Sciences) and Channel~06 \emph{Vidhik} (Law, Legal Studies \& Human Rights) are correctly matched. Channel~10 \emph{Daksh} is devoted to \emph{Applied Sciences}; Life Sciences is carried by Channel~09 \emph{Spandan}. Hence pair (D) is incorrect.
\tip{SWAYAM Prabha: 40 free 24$\times$7 educational DTH channels via GSAT-15, launched 9 July 2017; content from CEC--UGC, NPTEL (IITs), IGNOU, NCERT, NIOS and others.}

%% Q4
All three are plagiarism-detection tools used by Indian institutions:
\begin{itemize}
\item \textbf{Turnitin} -- widely used international service.
\item \textbf{DrillBit} -- Indian-developed software; provided free to universities by the Ministry of Education/INFLIBNET under the \emph{ShodhShuddhi} initiative (from 2022).
\item \textbf{Urkund} (now \emph{Ouriginal}) -- the tool originally supplied by INFLIBNET under ShodhShuddhi (2019).
\end{itemize}
The UGC (Promotion of Academic Integrity and Prevention of Plagiarism in HEIs) Regulations, 2018 require such checks for theses and dissertations. Hence I, II and III.

%% Q5
Storage units from the largest to the smallest:
\[\text{Terabyte (iv)} > \text{Gigabyte (v)} > \text{Megabyte (iii)} > \text{Kilobyte (i)} > \text{Byte (ii)} > \text{Bit (vi)}\]
(1~TB = 1024~GB, 1~GB = 1024~MB, 1~MB = 1024~KB, 1~KB = 1024~bytes, 1~byte = 8~bits.) The order iv, v, iii, i, ii, vi is option (A).

%% Q6
\textbf{Context} is the setting (physical, social, cultural, psychological and temporal) in which communication takes place. A sender who encodes a message keeping the circumstances and environment in mind is taking the \emph{context} into account; the same words may carry different meanings in different contexts.
\why{Decoding is the receiver interpreting the message; the channel is the medium that carries it; the response is the receiver's reaction or feedback.}

%% Q7
\textbf{Proxemics} (term coined by Edward T.~Hall, 1966) is the study of how people use \emph{space and distance} in communication: intimate, personal, social and public zones. Students keeping a distance to maintain their ``territory'' is a proxemic behaviour.
\why{Haptics = touch; Kinesics = body movements, gestures and facial expressions; Chronemics = use of time.}

%% Q8
Communication that flows from subordinates to superiors (teacher $\rightarrow$ head of institution) is \textbf{upward communication}. Reports, feedback, suggestions and grievances travel upward.
\why{Downward flows from superior to subordinate (circulars, instructions); horizontal/lateral flows between people at the same level (teacher to teacher).}

%% Q9
A \textbf{perceptual barrier} arises because different people look at the \emph{same} event, message or phenomenon from \emph{different frames of reference} (their experiences, beliefs, culture and expectations) and therefore interpret it differently.
\tip{Frame of reference = the mental ``lens'' formed by one's background. Same reality + different lenses = different perceptions.}

%% Q10
\textbf{Yellow journalism} is sensational, exaggerated and poorly researched reporting meant to attract readers; the term comes from the 1890s circulation war in New York between Pulitzer's \emph{New York World} and Hearst's \emph{New York Journal}. Its modern form is \textbf{tabloid journalism}.
\why{Development journalism reports on development issues; investigative journalism involves in-depth, fact-based inquiry; print journalism is simply the medium.}

%% Q11
Doctor, Nurse and Patient are all \emph{persons} found in a hospital; \textbf{Hospital} (item~2) is the \emph{place} itself. Hence item~2 is the odd one out, which is option (A).

%% Q12
In Bhatta Mimamsa and Advaita Vedanta, \textbf{Anupalabdhi} (non-perception or non-cognition) is the independent \emph{pramana} (means of valid knowledge) through which the \textbf{abhava} (non-existence or absence) of a thing is known. For example, ``there is no pot on the floor'' is known by the non-perception of a pot where it would have been perceived.
\why{Arthapatti is postulation or presumption (e.g.\ ``fat Devadatta does not eat by day, so he eats at night''). Upalabdhi means perception, which is the opposite. Upasthiti is not a pramana.}
\tip{Six pramanas (Mimamsa/Vedanta): Pratyaksha, Anumana, Upamana, Shabda, Arthapatti, Anupalabdhi.}

%% Q13
Fill the blanks with \textbf{s, q, p, q, r}:
\[\text{pqr\underline{s}}\ \ \text{\underline{q}rsp}\ \ \text{rs\underline{p}\underline{q}}\ \ \text{spq\underline{r}}\]
Each block of four letters is the previous block shifted one place to the left (pqrs $\rightarrow$ qrsp $\rightarrow$ rspq $\rightarrow$ spqr). The missing letters are s, q, p, q, r, i.e.\ ``sqpqr'', option (C).

%% Q14
The relation is \emph{thing : place where it is stored or kept}. Ideas are stored in the brain; money is kept (deposited) in a bank.
\why{Cash is a form of money, a loan is borrowed money, and an investment is a use of money. None of these is the place where money is kept.}

%% Q15
The \textbf{Nyaya five-membered syllogism (Panchavayava)} in its correct order:
\begin{enumerate}
\item \textbf{Pratijna} (proposition): The hill has fire.
\item \textbf{Hetu} (reason): Because it has smoke.
\item \textbf{Udaharana} (example): Wherever there is smoke there is fire, as in a kitchen.
\item \textbf{Upanaya} (application): The hill has smoke, which is invariably accompanied by fire.
\item \textbf{Nigamana} (conclusion): Therefore the hill has fire.
\end{enumerate}
This is option (B).

%% Q16
\textbf{Descriptive research} describes ``what is'', meaning the current status of a phenomenon as it naturally exists, without manipulating variables. Surveys, case studies, developmental and correlational studies are descriptive.
\why{Creating new theory is the aim of fundamental or grounded-theory research; proving laws under controlled conditions is experimental research; changing existing conditions is the aim of action research.}

%% Q17
An \textbf{opinionnaire} (attitude scale or opinion questionnaire) collects the opinions of a large number of respondents on social issues. Collecting such data from a sample to describe a population is the \textbf{survey method}, a form of descriptive research.
\why{The historical method studies past records; the experimental method manipulates variables; ``quantitative'' is a type of data or approach, not a specific method of this kind.}

%% Q18
Research ethics require \textbf{confidentiality and anonymity}: participants should be referred to by codes, pseudonyms or anonymous identifiers (e.g.\ Respondent R-12, Teacher T3) so that responses cannot be traced back to individuals.
\why{Listing names with responses violates confidentiality; asking participants to publish their own data is irrelevant; replacing names with those of the data collectors is misleading.}

%% Q19
Qualitative measurement deals with \emph{non-numerical} attributes. \textbf{Values, attitudes and anecdotes} are all qualitative, described in words, categories or ratings of quality.
\why{Every other set contains quantitative items: capital investment, profit, performance scores and expenditure are measured in numbers.}

%% Q20
The normal curve is symmetric about the mean (0~SD), so exactly half the observations lie on each side. The area from the mean to $-3\,\sigma$ is 49.87\%, which is broadly \textbf{50\%}.
\begin{center}
\begin{tikzpicture}
\begin{axis}[width=8cm,height=3.6cm,axis lines=middle,ytick=\empty,xtick={-3,-2,-1,0,1,2,3},
  xticklabels={$-3\sigma$,$-2\sigma$,$-1\sigma$,$0$,$1\sigma$,$2\sigma$,$3\sigma$},
  ymin=0,ymax=0.45,xmin=-3.5,xmax=3.5,tick label style={font=\scriptsize}]
\addplot[domain=-3:0,samples=60,fill=brand!25,draw=none] {exp(-x^2/2)/sqrt(2*pi)} \closedcycle;
\addplot[domain=-3.5:3.5,samples=100,thick,brand] {exp(-x^2/2)/sqrt(2*pi)};
\node[font=\scriptsize] at (axis cs:-1.1,0.12) {49.87\%};
\end{axis}
\end{tikzpicture}
\end{center}
\tip{Mean to $\pm1\sigma$: 34.13\% each side (68.26\% total); $\pm2\sigma$: 95.44\%; $\pm3\sigma$: 99.73\%.}

%% Q21
The \textbf{Gujarat Public Universities Bill, 2023} was passed by the Gujarat Legislative Assembly in \textbf{September 2023} and became the \emph{Gujarat Public Universities Act, 2023}. It brings 11 state public universities (e.g.\ Gujarat University, M.S.~University of Baroda, Saurashtra University) under a common Act, replacing their separate Acts, with a Board of Management and a Chancellor appointed by the State.

%% Q22
NEP~2020 states its vision as an India-centred education system that contributes directly to \textbf{transforming India into an equitable and vibrant knowledge society}, by providing high-quality education to all and making India a global knowledge superpower.
\why{It does not seek an expensive system; ``free education for all'' is not its primary focus (RTE already covers ages 6--14); it blends traditional knowledge with modern education rather than emphasising only tradition.}

%% Q23
SWAYAM Prabha channels are produced by partner institutions. Channels coordinated by CEC--UGC are prepared through its \textbf{Educational Multimedia Research Centres (EMRCs)}. Channel~7 (\emph{Kautilya}) is handled by the \textbf{EMRC, Gujarat University, Ahmedabad}.
\tip{CEC (Consortium for Educational Communication) is an inter-university centre of UGC with a network of about 21 EMRCs/AVRCs across India.}

%% Q24
\textbf{IGNOU}, the Indira Gandhi National Open University (est.~1985, New Delhi), is a central open university and the largest in the world by enrolment. It offers the largest number of programmes (over 200) through Open and Distance Learning. It also serves as the national apex body for ODL.
\why{Mumbai and Delhi universities offer distance courses through their distance-education wings, but far fewer. BAOU is Gujarat's state open university.}

%% Q25
\textbf{NEET-UG} (National Eligibility-cum-Entrance Test), conducted by the \textbf{National Testing Agency (NTA)}, is the single common entrance examination for MBBS, BDS and AYUSH courses in India.
\why{JEE is for engineering (IITs/NITs); CAT is for management (IIMs); GATE is for postgraduate engineering admissions and PSU recruitment.}

%% Q26
Year-to-year changes in sales:
\begin{center}\small
\begin{tabular}{lccc}\toprule
Period & From $\rightarrow$ To & Change & \% change\\\midrule
2017--18 & 48 $\rightarrow$ 40 & $-8$ & $-16.7\%$\\
2018--19 & 40 $\rightarrow$ 30 & $-10$ & $-25.0\%$\\
2019--20 & 30 $\rightarrow$ 25 & $-5$ & $-16.7\%$\\
\textbf{2020--21} & \textbf{25 $\rightarrow$ 28} & $\mathbf{+3}$ & $\mathbf{+12.0\%}$\\
2021--22 & 28 $\rightarrow$ 40 & $+12$ & $+42.9\%$\\\bottomrule
\end{tabular}
\end{center}
The smallest change, in both absolute and percentage terms, is between \textbf{2020 and 2021}.

%% Q27 fig
\begin{tabular}{ccc}\toprule
\textbf{Location} & \textbf{Candidates (lakh)} & \textbf{Pass : Fail}\\\midrule
P & 1.25 & 7 : 3\\ Q & 3.14 & 5 : 3\\ R & 1.08 & 4 : 5\\ S & 2.27 & 1 : 3\\ T & 1.85 & 3 : 2\\\bottomrule
\end{tabular}

%% Q27 after
What is the ratio of the number of candidates failing the exam from location S to that of those failing the exam from location P?

%% Q27
Failed from S $= 2.27 \times \dfrac{3}{1+3} = 2.27\times\dfrac34 = 1.7025$ lakh.\\
Failed from P $= 1.25 \times \dfrac{3}{7+3} = 1.25\times\dfrac{3}{10} = 0.375$ lakh.
\[\frac{1.7025}{0.375} = \frac{2.27\times 3/4}{1.25\times 3/10} = \frac{2.27\times 10}{1.25\times 4} = \frac{22.7}{5} = \frac{227}{50}\]
Ratio S : P $=$ \textbf{227 : 50}.

%% Q28
This is Bayes' theorem. Let $C$ = correct diagnosis and $D$ = death.
\[P(C)=0.6,\quad P(D\,|\,C)=0.4;\qquad P(\bar C)=0.4,\quad P(D\,|\,\bar C)=0.7\]
\[P(C\,|\,D)=\frac{0.6\times0.4}{0.6\times0.4+0.4\times0.7}=\frac{0.24}{0.24+0.28}=\frac{0.24}{0.52}\approx \mathbf{0.46}\]
\begin{center}
\begin{tikzpicture}[font=\scriptsize,>=Stealth,node distance=4mm]
\node (p) {Patient with X};
\node (c) at (3,0.6) {Correct (0.6)}; \node (w) at (3,-0.6) {Wrong (0.4)};
\node (cd) at (6.2,0.6) {Dies (0.4): $0.24$}; \node (wd) at (6.2,-0.6) {Dies (0.7): $0.28$};
\draw[->] (p.east) -- (c.west); \draw[->] (p.east) -- (w.west);
\draw[->] (c) -- (cd); \draw[->] (w) -- (wd);
\end{tikzpicture}
\end{center}

%% Q29
$P(\text{not swimmer})=\tfrac15 \Rightarrow p=P(\text{swimmer})=\tfrac45=0.8,\ q=0.2$, $n=5$ (binomial).
\[P(X\ge4)=\binom54(0.8)^4(0.2)+\binom55(0.8)^5=5(0.4096)(0.2)+0.32768=0.4096+0.32768=0.73728\]
$\approx$ \textbf{0.74}.

%% Q30
Multiples of 3 in 1--100: $\lfloor100/3\rfloor=33$. Multiples of 5: $20$. Multiples of both (i.e.\ of 15): $6$.
\[n(3\cup5)=33+20-6=47\qquad\Rightarrow\qquad P=\frac{47}{100}=\mathbf{0.47}\]

%% Q31
Memory is studied through the processes of \textbf{encoding/learning, retention (storage), recall and recognition}, and these are the recognised states or levels at which memory is measured. \textbf{Assimilation} is a Piagetian concept of \emph{cognitive development} (fitting new experience into existing schemas). It is not a state of memory.

%% Q32
Recognising individual differences (in ability, pace, learning style, interest and background) lets a teacher \textbf{guide learning} appropriately and \textbf{develop each student's cognitive processes} through differentiated instruction.
\why{Classifying students into subgroups is at best a tool, not the purpose; uniform teaching ignores differences; individual differences do not reduce the need for assessment.}

%% Q33
Modern learning theory (constructivism) holds that learning happens when learners \textbf{actively} construct knowledge by doing, discussing, questioning and reflecting. Active learning processes therefore matter most, while background, prior achievement and textbooks are supporting factors.

%% Q34
SWAYAM Prabha telecasts high-quality \textbf{educational programmes in diverse disciplines} (arts, science, commerce, engineering, social sciences, law, medicine, agriculture and more) for school, higher-education and lifelong learners. For the higher-education channels the target group is learners aged roughly 18 and above.

%% Q35
Under the \textbf{Choice Based Credit System (CBCS)}, introduced by UGC in 2015, each course carries credits, and a degree is awarded when the student earns the \textbf{prescribed minimum number of credits} across core, elective and ability- or skill-enhancement courses. Grades are expressed as SGPA/CGPA.
\why{Years of study, raw marks or attendance alone do not decide the award of the degree under CBCS.}

%% Q36
Fly ash, the fine residue of coal combustion, is used for:
\begin{itemize}
\item \textbf{road and embankment construction} (the MoRTH/IRC encourage its use);
\item \textbf{wasteland reclamation and development};
\item \textbf{agriculture}, as a soil conditioner supplying micro-nutrients (B, Mo, Zn, S) and improving water-holding capacity;
\item \textbf{brick making}: fly-ash bricks are mandated near thermal power plants by the MoEF notification of 1999.
\end{itemize}
All four uses are correct, so the answer is (C).

%% Q37
SDG numbers: Life on land = \textbf{15}, Life below water = \textbf{14}, Climate action = \textbf{13}, Affordable and clean energy = \textbf{7}, Clean water and sanitation = \textbf{6}.
Decreasing order: 15 (III), 14 (I), 13 (IV), 7 (V), 6 (II), giving \textbf{III, I, IV, V, II}.
\tip{SDGs 6--7--13--14--15 = Water, Energy, Climate, Oceans, Land.}

%% Q38
\textbf{Statement I} is true: cyclones that originate over warm tropical oceans, in the belt between the Tropic of Cancer and the Tropic of Capricorn, are called tropical cyclones (hurricanes, typhoons and so on).
\textbf{Statement II} is true: Hurricane Helene (late September 2024) made landfall in Florida's Big Bend as a Category~4 storm and caused catastrophic flooding and more than 200 deaths across the south-eastern US, especially North Carolina.
Both statements are true, so the answer is (A).

%% Q39
The Department of Science \& Technology's national report \emph{Climate Vulnerability Assessment for Adaptation Planning in India using a Common Framework} identified \textbf{eight} eastern and north-eastern states as highly vulnerable: Jharkhand, Mizoram, Odisha, Chhattisgarh, Assam, Bihar, Arunachal Pradesh and West Bengal.

%% Q40 opts
Nitric Oxide (NO)
Nitrous Oxide (N$_2$O)
Nitrogen Dioxide (NO$_2$)
Dinitrogen Tetroxide (N$_2$O$_4$)

%% Q40
India's \textbf{National Air Quality Index} (launched 2014) uses eight pollutants: PM$_{10}$, PM$_{2.5}$, \textbf{NO$_2$}, SO$_2$, CO, O$_3$, NH$_3$ and Pb. Of the nitrogen oxides, only \textbf{nitrogen dioxide} is an AQI constituent.
\tip{AQI categories: Good (0--50), Satisfactory (51--100), Moderately polluted (101--200), Poor (201--300), Very poor (301--400), Severe (401--500).}

%% Q41
COMMUNICATIONS has 14 letters. Swap each pair:
\begin{center}\ttfamily
\begin{tabular}{*{14}{c}}
C&O&M&M&U&N&I&C&A&T&I&O&N&S\\
O&C&M&M&N&U&C&I&T&A&O&I&S&N\\
\end{tabular}
\end{center}
10th from the right = $(14-10+1)$ = 5th from the left = \textbf{N}.

%% Q42
$A=B+C$ and $A+B+C=180^\circ \Rightarrow 2A=180^\circ \Rightarrow A=90^\circ$, so $B+C=90^\circ$.
With $B:C=4:5$: $B=\tfrac49\times90^\circ=40^\circ$, $C=\tfrac59\times90^\circ=50^\circ$. The angles are \textbf{90\textdegree, 40\textdegree, 50\textdegree}.

%% Q43
The series alternately multiplies by 2 and by 1.5:
\[4\xrightarrow{\times2}8\xrightarrow{\times1.5}12\xrightarrow{\times2}24\xrightarrow{\times1.5}36\xrightarrow{\times2}72\xrightarrow{\times1.5}\mathbf{108}\]
(Equivalently, the differences run $+4,+4,+12,+12,+36,+36$.)

%% Q44
Original saving $=13{,}500-9{,}000=4{,}500$.\\
New income $=13{,}500\times1.14=15{,}390$; new expenditure $=9{,}000\times1.07=9{,}630$.\\
New saving $=15{,}390-9{,}630=5{,}760$.
\[\text{Increase}=\frac{5{,}760-4{,}500}{4{,}500}\times100=\frac{1{,}260}{4{,}500}\times100=\mathbf{28\%}\]

%% Q45
Half-yearly compounding: rate $=10/2=5\%$ per half-year, periods $=18$ months $\div 6 = 3$.
\[A=1000(1.05)^3=1000\times1.157625=1157.63\qquad CI=1157.63-1000=\textbf{Rs.\ 157.63}\]

%% Q46 ctx
\textbf{Directions for Q.~46 to 50:} Read the following paragraph and answer the questions that follow.\par\smallskip
Effective learning requires a comprehensive approach that involves appropriate curriculum, engaging pedagogy, continuous formative assessment, and adequate student support. The curriculum must be interesting and relevant, and updated regularly to align with the latest knowledge requirements and to meet specified learning outcomes. High-quality pedagogy is then necessary to successfully impart the curricular material to students; pedagogical practices determine the learning experiences that are provided to students, thus directly influencing learning outcomes. The assessment methods must be scientific, designed to continuously improve learning and test the application of knowledge.

To promote creativity, institutions and faculty will have the autonomy to innovate on matters of curriculum, pedagogy, and assessment within a broad framework of higher education qualifications that ensures consistency across institutions and programmes and across the ODL, online, and traditional `in-class' modes. All assessment systems shall also be decided by the HEI, including those that lead to final certification. The Choice Based Credit System (CBCS) will be revised for instilling innovation and flexibility. HEIs shall move to a criterion-based grading system that assesses student's achievement based on the learning goals for each programme, making the system fairer and outcomes more comparable.

%% Q46
The opening sentence lists the four components: ``appropriate \textbf{curriculum}, engaging \textbf{pedagogy}, continuous formative \textbf{assessment}, and adequate \textbf{student support}.''
(This passage is taken from NEP~2020, paragraph~12.1.)

%% Q47
The passage says pedagogical practices ``determine the \textbf{learning experiences} that are provided to students, thus directly influencing \textbf{learning outcomes}.''
\why{Grades, homework and classroom interaction are not what the passage names.}

%% Q48
``To promote creativity, institutions and faculty will have the autonomy to innovate on matters of \textbf{curriculum, pedagogy, and assessment}.''

%% Q49
``The Choice Based Credit System (CBCS) will be revised for instilling \textbf{innovation} and flexibility.''

%% Q50
``HEIs shall move to a \textbf{criterion-based grading system} that assesses student's achievement based on the learning goals for each programme.'' Criterion-referenced grading measures attainment against fixed learning goals rather than against other students (norm-referenced grading).
